% ======================================================================
\chapter{Padrões de commit}
\label{cap:padroes}

No capítulo~\ref{cap:commits} vimos como fazer commits. Este capítulo trata de \textbf{como escrever essa mensagem} de um jeito que a equipe inteira (e as ferramentas) consigam ler.

\section{Por que padronizar}

Compare dois históricos do mesmo projeto:

\begin{saida}
a41aded ajustes
c2524a8 arrumei
afe16bb agora vai
4203776 mudanças do joão
\end{saida}
\begin{saida}
a41aded fix(login): corrige senha vazia
c2524a8 feat(relatorio): exporta em PDF
afe16bb docs: explica a instalação
4203776 test(login): cobre senha vazia
\end{saida}

No segundo, dá para saber \textbf{o tipo} da mudança, \textbf{onde} ela aconteceu e \textbf{o que} foi feito, sem abrir um arquivo sequer. Um padrão de commit também permite automatizar tarefas: gerar notas de versão, decidir o próximo número de versão e barrar mensagens fora do combinado.

\begin{conceito}[A regra mais importante]
O melhor padrão é \textbf{o que a equipe combinou e segue}. Os padrões deste capítulo são os mais usados, mas consistência vale mais que qualquer escolha isolada.
\end{conceito}

\section{Conventional Commits}

O padrão mais difundido é o \textbf{Conventional Commits} 1.0.0 (\url{https://www.conventionalcommits.org/pt-br/v1.0.0/}). A estrutura é:

\begin{bashcode}
<tipo>[escopo opcional]: <descrição>

[corpo opcional]

[rodapé(s) opcional(is)]
\end{bashcode}

As regras principais da especificação:

\begin{itemize}
  \item o \textbf{tipo} é obrigatório e vem primeiro, seguido de \textbf{dois-pontos e espaço};
  \item o \textbf{escopo} é opcional: um substantivo entre parênteses que diz qual parte do sistema mudou, como \arq{fix(parser):};
  \item a \textbf{descrição} vem logo depois de \arq{: } e resume a mudança;
  \item o \textbf{corpo} é opcional e começa \textbf{uma linha em branco} depois da descrição, com quantos parágrafos forem necessários;
  \item os \textbf{rodapés} são opcionais e vêm uma linha em branco depois do corpo, no formato \arq{Token: valor} (por exemplo, \arq{Refs: \#123} ou \arq{Reviewed-by: Maria});
  \item uma \textbf{quebra de compatibilidade} (\emph{breaking change}) é indicada com \arq{!} antes dos dois-pontos, com um rodapé \arq{BREAKING CHANGE: <explicação>}, ou com os dois.
\end{itemize}

\begin{bashcode}
feat(api)!: exige autenticação em todos os endpoints

Antes, os endpoints de consulta eram públicos. Agora todo pedido
precisa do cabeçalho Authorization com um token válido.

BREAKING CHANGE: clientes sem token passam a receber 401.
Refs: #42
\end{bashcode}

\subsection{Os tipos}

A especificação só define dois tipos: \textbf{\arq{feat}} (novo recurso) e \textbf{\arq{fix}} (correção de bug). Os outros são convenção, herdada das regras do projeto Angular. A lista abaixo é a aceita pela configuração padrão do \emph{commitlint} (seção \ref{sec:commitlint}):

\begin{center}
\renewcommand{\arraystretch}{1.3}
\begin{tabularx}{\linewidth}{@{}>{\ttfamily\color{roxoescuro}}l X@{}}
\toprule
\textrm{\textbf{Tipo}} & \textbf{Quando usar} \\
\midrule
feat & um recurso novo para quem usa o sistema \\
fix & a correção de um bug \\
docs & só documentação (README, comentários, manuais) \\
style & formatação que não muda o comportamento (espaços, ponto e vírgula, indentação) \\
refactor & reorganização do código sem mudar o que ele faz nem corrigir bug \\
perf & melhoria de desempenho \\
test & criação ou ajuste de testes \\
build & sistema de build ou dependências (npm, pip, Maven\ldots) \\
ci & configuração de integração contínua (GitHub Actions) \\
chore & tarefas de manutenção que não se encaixam nos outros tipos \\
revert & desfaz um commit anterior \\
\bottomrule
\end{tabularx}
\end{center}

\begin{dica}[Tipos e versões]
Com Conventional Commits, ferramentas conseguem calcular a próxima versão pelo \emph{Versionamento Semântico}: \arq{fix} sobe o número de correção (1.2.\textbf{3}), \arq{feat} sobe o número de recurso (1.\textbf{3}.0) e um \arq{BREAKING CHANGE} sobe o número principal (\textbf{2}.0.0).
\end{dica}

\begin{aviso}[Tipos fora da lista]
Algumas referências populares em português acrescentam tipos como \arq{raw}, \arq{cleanup}, \arq{remove} e \arq{init}. Eles não fazem parte da especificação nem da configuração padrão do commitlint, que os rejeita. Só use se a equipe combinar e configurar isso.
\end{aviso}

\section{Boas práticas para a mensagem}

\begin{itemize}
  \item \textbf{Uma mudança por commit.} Se a descrição precisa de um \enquote{e}, provavelmente são dois commits. Commits pequenos são mais fáceis de revisar e de desfazer.
  \item \textbf{Descrição curta.} O livro \emph{Pro Git} recomenda uma primeira linha de até uns 50 caracteres e o corpo quebrado em 72 colunas. A configuração padrão do commitlint aceita até 100 caracteres no cabeçalho.
  \item \textbf{O corpo explica o porquê.} O \enquote{o quê} está no \comando{git diff}. O motivo da mudança só existe na mensagem.
  \item \textbf{Um único modo verbal.} As convenções em inglês usam o imperativo (\enquote{add}, \enquote{fix}). Em português, equipes usam o imperativo (\enquote{corrige}), o presente (\enquote{corrige}) ou o gerúndio (\enquote{corrigindo}). A especificação não exige nenhum. Escolha um e mantenha.
  \item \textbf{Idioma combinado.} A especificação também não fixa idioma. Misturar português e inglês no mesmo histórico atrapalha a busca.
\end{itemize}

\section{Gitmoji: emojis nas mensagens}

O \textbf{gitmoji} (\url{https://gitmoji.dev}) associa um emoji a cada tipo de mudança. No texto do commit, o emoji aparece como um código entre dois-pontos, o \emph{shortcode}. Alguns dos mais usados, com a descrição oficial traduzida:

\begin{center}
\renewcommand{\arraystretch}{1.25}
\small
\begin{tabularx}{\linewidth}{@{}>{\ttfamily\color{roxoescuro}}l X >{\ttfamily\color{roxoescuro}}l X@{}}
\toprule
\textrm{\textbf{Código}} & \textbf{Significado} & \textrm{\textbf{Código}} & \textbf{Significado} \\
\midrule
:sparkles: & novo recurso & :bug: & corrige um bug \\
:memo: & documentação & :art: & estrutura ou formato do código \\
:zap: & desempenho & :recycle: & refatoração \\
:fire: & remove código ou arquivos & :truck: & move ou renomeia arquivos \\
:white\_check\_mark: & testes & :construction\_worker: & CI \\
:wrench: & arquivos de configuração & :package: & arquivos compilados ou pacotes \\
:lock: & segurança ou privacidade & :ambulance: & correção crítica urgente \\
:pencil2: & corrige erro de digitação & :rewind: & reverte mudanças \\
:tada: & começa um projeto & :bookmark: & versão ou tag \\
:page\_facing\_up: & licença & :card\_file\_box: & mudanças de banco de dados \\
\bottomrule
\end{tabularx}
\end{center}

A lista completa tem 75 emojis e está em \url{https://gitmoji.dev}.

\begin{aviso}[Três cuidados com emojis]
\begin{enumerate}
  \item \textbf{O terminal não converte.} O GitHub mostra o emoji na página do commit, mas o \comando{git log} exibe o texto \arq{:sparkles:} como foi escrito.
  \item \textbf{O mesmo emoji muda de sentido entre fontes.} Na lista oficial, \arq{:boom:} é quebra de compatibilidade, mas há tutoriais que o usam para \enquote{reverter}. Siga a lista oficial.
  \item \textbf{Emoji antes do tipo quebra o padrão.} A especificação exige que a mensagem comece pelo tipo. Uma mensagem como \arq{:bug: fix: corrige o login} é rejeitada pelo commitlint na configuração padrão.
\end{enumerate}
\end{aviso}

\section{Um padrão combinando os dois}

Muitas equipes juntam gitmoji e Conventional Commits num padrão próprio. É o caso dos repositórios do autor deste guia, que usam:

\begin{bashcode}
<:emoji:> <tipo>(<escopo>): <Descrição no gerúndio>

:memo: docs(guia): Adicionando os exercícios de assinatura
:bug: fix(login): Corrigindo o acesso com senha vazia
:truck: chore(docs): Movendo o manual para a pasta materials
\end{bashcode}

Não há nada de errado nisso, desde que \textbf{seja uma decisão da equipe, esteja escrita} (no README ou num guia de contribuição) e \textbf{as ferramentas estejam configuradas para aceitá-la}. O ponto é a consistência.

\section{Ferramentas}
\label{sec:commitlint}

\subsection*{commitlint: barrar mensagens fora do padrão}

O \textbf{commitlint} (\url{https://commitlint.js.org}) confere cada mensagem antes de o commit ser gravado. Num projeto com Node.js:

\begin{bashcode}
npm install -D @commitlint/cli @commitlint/config-conventional
echo "export default { extends: ['@commitlint/config-conventional'] };" > commitlint.config.js

# testar uma mensagem sem precisar fazer commit
echo "feat: adiciona exportação em PDF" | npx commitlint
\end{bashcode}

Para rodar automaticamente em todo \comando{git commit}, o guia oficial usa o \emph{husky}:

\begin{bashcode}
npm install --save-dev husky
npx husky init
echo "npx --no -- commitlint --edit \$1" > .husky/commit-msg
\end{bashcode}

Depois disso, \comando{git commit -m "foo: teste"} é recusado, porque \arq{foo} não é um tipo válido.

\begin{aviso}[A configuração padrão é rígida]
Além de exigir o tipo no início, a configuração padrão rejeita descrições que começam com letra maiúscula. Um padrão como o da seção anterior precisa de regras próprias no \arq{commitlint.config.js}.
\end{aviso}

\subsection*{gitmoji-cli: escolher o emoji num menu}

\begin{bashcode}
npm i -g gitmoji-cli
gitmoji -c        # monta o commit de forma interativa
gitmoji -l        # lista os emojis
\end{bashcode}

\subsection*{Sem ferramenta nenhuma}

Nada disso é obrigatório. Um arquivo \arq{CONTRIBUTING.md} com o padrão e três exemplos já resolve boa parte do problema, e a revisão do Pull Request (capítulo~\ref{cap:commits}) é um bom momento para lembrar o combinado.
